\label{sec:theory}

We now provide rigorous theoretical analysis characterizing when and why edge-weighting produces method-independent solutions. We formalize the conditions under which all partitioning algorithms converge 	o identical outcomes, derive the phase transition point, and analyze computational implications.

\subsection{Method Convergence: Formal Characterization}

We begin by formalizing the phenomenon observed empirically: for sufficiently large edge weight factor $\alpha$, all partitioning methods produce identical results.

\begin{definition}[Weighted Graph Partitioning Instance]
A weighted graph partitioning instance is a tuple $(G, w, k, \epsilon)$ where:
\begin{itemize}
    \item $G = (V, E)$ is an undirected graph with $|V| = n$ vertices
    \item $w: E \to \mathbb{R}^+$ assigns positive weights 	o edges
    \item $k$ is the number of partitions (districts)
    \item $\epsilon > 0$ is the balance constraint tolerance
\end{itemize}
\end{definition}

For VRA-compliant redistricting, we consider a specific weight structure:
\[
w(u,v) = \begin{cases}
\alpha & \text{if } m_u \geq \tau \text{ and } m_v \geq \tau \\
1 & \text{otherwise}
\end{cases}
\]
where $m_v \in [0,1]$ is the minority fraction at vertex $v$, $\tau$ is the minority threshold (e.g., 0.40), and $\alpha \gg 1$ is the edge weight factor.

\begin{theorem}[Method Equivalence Conditions]
\label{thm:convergence}
Let $(G, w, k, \epsilon)$ be a weighted partitioning instance with edge weights as above. Suppose the following structural conditions hold:
\begin{enumerate}[label=\textbf{(C\arabic*)}]
    \item \textbf{(Connectivity)}: The full graph $G$ is connected and the minority subgraph $G_m = (V, E_m)$ has exactly $q \geq 1$ connected components $S_1,\ldots,S_q$, each spatially contiguous.
    \item \textbf{(Separation)}: Every minority component $S_i$ satisfies $|E(S_i, V \setminus S_i) \cap E_m| = 0$, i.e., no minority--minority edge crosses component boundaries (components are edge-separated in $G_m$).
    \item \textbf{(Balance feasibility)}: There exists a partition $P^*$ into $k$ parts respecting the population balance constraint $|N_i - N/k| \leq \epsilon N/k$ such that each minority component $S_i$ lies entirely within a single part of $P^*$.
\end{enumerate}
Define the \textbf{weight ratio} $\rho = \alpha$ and the \textbf{critical threshold}
\[
    \alpha_{\text{crit}}(G,k,\epsilon) \;=\; \frac{|E_r^*|}{\epsilon \cdot |E_m^*| / k}
\]
where $|E_r^*|$ is the number of regular boundary edges in $P^*$ and $|E_m^*|$ is the minimum number of minority--minority edges that any balance-feasible partition must cut to achieve VRA-required minority concentration. Under conditions \textup{(C1)--(C3)}, for $\alpha \geq \alpha_{\text{crit}}$:

\begin{enumerate}
    \item \textbf{(Objective Convergence)}: Every algorithm $\mathcal{A}$ satisfying the balance constraint produces a partition $P_\mathcal{A}$ with objective value
    \[
    \mathrm{Obj}(P_{\mathcal{A}}) \;\leq\; \mathrm{Obj}^* + \delta(\alpha),
    \qquad \delta(\alpha) = O\!\left(\frac{|E_r|}{\alpha}\right) = O\!\left(\frac{1}{\alpha}\right),
    \]
    where $\mathrm{Obj}^*$ is the minimum weighted cut over all balance-feasible partitions.

    \item \textbf{(Partition Uniqueness)}: If condition \textup{(C3)} holds with a \emph{unique} such feasible assignment (i.e., the minority components $S_1,\ldots,S_q$ can be assigned to districts in only one way while satisfying population balance), then $P_{\mathcal{A}_1} = P_{\mathcal{A}_2}$ for any two algorithms $\mathcal{A}_1, \mathcal{A}_2$ both satisfying the balance constraint.
\end{enumerate}
\end{theorem}

\begin{proof}[Proof Sketch]
\textbf{Part 1 (Objective Convergence)}.

Let $P$ be any balance-feasible partition. Decompose its cut into minority--minority edges and regular edges:
\[
    \mathrm{Obj}(P) = \alpha \cdot c_m(P) + c_r(P),
\]
where $c_m(P) = |\{(u,v)\in E_m : u,v \text{ in different parts of } P\}|$ and $c_r(P)$ counts regular boundary edges.

Under condition (C3), the optimal partition $P^*$ has $c_m(P^*) = 0$ (it places every minority component in a single district), so $\mathrm{Obj}^* = c_r(P^*)$. Any competing partition $P$ with $c_m(P) \geq 1$ incurs
\[
    \mathrm{Obj}(P) \geq \alpha + c_r(P) \geq \alpha.
\]
For $P^*$, a balance-maintaining perturbation that cuts $j$ additional regular edges while avoiding all minority edges costs $\mathrm{Obj}^* + j$. We need to show that any algorithm produces a partition within $\delta(\alpha)$ of $P^*$.

Since $\mathrm{Obj}^* = c_r(P^*) \leq |E_r|$ and the cost of cutting any minority edge is $\alpha \gg \mathrm{Obj}^*$ for $\alpha \geq \alpha_{\text{crit}}$, any local-search algorithm (including METIS's Kernighan--Lin refinement) that moves vertices to reduce the weighted cut will avoid any move that cuts a minority edge. The algorithm is therefore constrained to solutions with $c_m = 0$, and the only flexibility lies in the assignment of non-minority vertices to achieve population balance. The cost of this assignment is at most $\mathrm{Obj}^* + O(|E_r|/(\alpha\epsilon)) = \mathrm{Obj}^* + O(1/\alpha)$, giving $\delta(\alpha) = O(|E_r|/\alpha)$.

\textbf{Part 2 (Partition Uniqueness)}.

When the assignment of minority components to districts is unique (condition C3 with uniqueness), there is exactly one partition $P^*$ with $c_m = 0$ that satisfies population balance. Since every algorithm producing a balance-feasible partition with $\alpha \geq \alpha_{\text{crit}}$ must have $c_m = 0$ (by Part 1, any partition with $c_m \geq 1$ has objective at least $\alpha > \mathrm{Obj}^* + \delta(\alpha)$, contradicting convergence), every such algorithm must produce $P^*$. Hence $P_{\mathcal{A}_1} = P_{\mathcal{A}_2} = P^*$.
\end{proof}

\begin{remark}[Role of Conditions (C1)--(C3)]
\label{rem:conditions}
Condition (C1) ensures the weighted Laplacian's Fiedler value grows with $\alpha$ (Proposition~\ref{prop:spectral}). Condition (C2) ensures that separate minority clusters are ``eigenvalue-separated'' in $L_w$, guaranteeing the large gap $\lambda_2/\lambda_3 = \Theta(\alpha)$. Condition (C3) is the decisive balance-feasibility condition: it asserts that minority concentration can be achieved without cutting any minority--minority edge. When (C3) holds, edge-weighting dominates algorithmic choice regardless of whether one uses recursive bisection, adaptive tree selection, or $k$-way optimization---all three are forced to find the same cut. When (C3) fails (e.g., minority population is too dispersed to form VRA districts without splitting a minority cluster), edge-weighting alone cannot guarantee method equivalence and an $\alpha$-sweep may reveal residual method variance.
\end{remark}

\subsection{Phase Transition: Characterizing $\alpha_{\text{crit}}$}

Theorem~\ref{thm:convergence} establishes existence of $\alpha_{\text{crit}}$ but does not provide a formula. We now derive bounds on the critical threshold.

\begin{theorem}[Phase Transition]
\label{thm:phase-transition}
Define the variance in maximum minority percentage across algorithms as:
\[
V(\alpha) = \text{Var}_{\mathcal{A} \in \text{Algorithms}} \left[ \max_{i \in [k]} \frac{M_i(P_{\mathcal{A}})}{N_i(P_{\mathcal{A}})} \right]
\]
where $M_i, N_i$ are minority population and total population in district $i$.

There exists a sharp phase transition:
\[
V(\alpha) = \begin{cases}
\Theta(1) & \text{if } \alpha < \alpha_{\text{crit}} \\
O(1/\alpha^2) & \text{if } \alpha > \alpha_{\text{crit}}
\end{cases}
\]

Furthermore, $\alpha_{\text{crit}}$ can be bounded as:
\[
\frac{k}{m_{\text{state}}} \leq \alpha_{\text{crit}} \leq \frac{k \cdot \Delta}{\epsilon \cdot m_{\text{state}}}
\]
where $m_{\text{state}}$ is the state-wide minority fraction, $\Delta$ is the maximum vertex degree, and $\epsilon$ is the balance tolerance.
\end{theorem}

\begin{proof}[Proof Sketch]
The lower bound follows from counting argument: to create $\frac{k \cdot m_{\text{state}}}{\tau}$ districts with minority concentration $\tau > m_{\text{state}}$, we must concentrate minority population. This requires $\alpha$ large enough that cutting minority edges is avoided. Since each district has $\frac{1}{k}$ of total population, and we want $\tau \approx 0.6$ while $m_{\text{state}} \approx 0.4$, we need:
\[
\alpha \geq \frac{k}{m_{\text{state}}}
\]

The upper bound accounts for graph structure: vertices have bounded degree $\Delta$, and balance tolerance $\epsilon$ allows some flexibility. Detailed analysis (see Appendix A) yields the multiplicative factor.

The phase transition is sharp because once $\alpha$ exceeds the threshold, the cost of cutting minority edges dominates all other considerations exponentially. This creates a discrete ``flip'' from method-dependent 	o method-independent regimes.
\end{proof}

\begin{corollary}[Practical Prediction]
For states with $m_{\text{state}} \approx 0.37$ (average across our test states) and $k \in [4, 14]$:
\[
\alpha_{\text{crit}} \in \left[ \frac{4}{0.37}, \frac{14}{0.37} \right] = [11, 38]
\]

However, empirical results show $\alpha_{\text{crit}} \approx 3-5$, suggesting the upper bound is loose. The tighter bound accounts for spatial clustering: minority-dense tracts form contiguous regions, reducing the effective degree $\Delta$ and tightening the transition point.
\end{corollary}

\subsection{Spectral Analysis: Eigenvalue Structure}
\label{sec:spectral}

We provide rigorous theoretical insight through spectral graph theory, connecting method convergence to the eigenvalue structure of the weighted Laplacian. The key result is that for sufficiently large $\alpha$, a large eigenvalue gap emerges that forces the Fiedler vector to encode the minority-dense subgraph structure, making the optimal balanced cut independent of the bisection tree used to find it.

\subsubsection{Setup and Notation}

Let $G = (V, E)$ be a connected graph with $n$ vertices. Partition the edges into $E = E_m \cup E_r$ where $E_m$ is the set of \emph{minority--minority} edges (both endpoints have minority fraction $\geq \tau$) and $E_r = E \setminus E_m$ is the remaining edge set. Define the weight function
\[
    w(u,v) = \begin{cases} \alpha & (u,v) \in E_m, \\ 1 & (u,v) \in E_r. \end{cases}
\]
Let $W$ be the $n\times n$ weighted adjacency matrix, $D_w = \mathrm{diag}(d_1,\ldots,d_n)$ the weighted degree matrix with $d_i = \sum_{j:(i,j)\in E} w(i,j)$, and $L_w = D_w - W$ the \emph{weighted Laplacian}. Write $L_w = L_r + (\alpha-1) L_m$ where $L_r$ and $L_m$ are the unweighted Laplacians of $G_r = (V,E_r)$ and $G_m = (V,E_m)$ respectively.

Denote by $0 = \lambda_1(L_w) \leq \lambda_2(L_w) \leq \cdots \leq \lambda_n(L_w)$ the eigenvalues of $L_w$, and by $\mathbf{f}_2 \in \mathbb{R}^n$ (normalized, $\|\mathbf{f}_2\|=1$) the corresponding Fiedler vector.

\subsubsection{Eigenvalue Gap Proposition}

\begin{proposition}[Eigenvalue Gap Under Edge Weighting]
\label{prop:spectral}
Let $G = (V,E)$ be a connected graph with the weight structure above, and assume the minority subgraph $G_m = (V,E_m)$ has $c \geq 1$ connected components with vertex sets $S_1,\ldots,S_c$ (vertices in $V \setminus (S_1\cup\cdots\cup S_c)$ are isolated in $G_m$). Let $\phi_m$ denote the edge expansion of $G_m$ restricted to its largest component, and let $\lambda_2^{(r)}$ denote the Fiedler value of $G_r$. Then:

\begin{enumerate}
    \item \textbf{(Lower bound on $\lambda_2$)}: There exists a constant $\beta > 0$ depending only on $G_m$'s component structure such that
    \[
        \lambda_2(L_w) \;\geq\; \beta\,\alpha - C_0
    \]
    for some constant $C_0 \geq 0$ independent of $\alpha$. In particular, $\lambda_2(L_w) = \Theta(\alpha)$.

    \item \textbf{(Upper bound on $\lambda_3$)}: If $G_m$ has at least two components (i.e., $c \geq 2$), then
    \[
        \lambda_3(L_w) \;\leq\; \lambda_3(L_r) + O(1),
    \]
    so $\lambda_3(L_w) = O(1)$ and the eigenvalue \emph{gap} satisfies
    \[
        \frac{\lambda_2(L_w)}{\lambda_3(L_w)} \;\xrightarrow{\alpha\to\infty}\; \infty.
    \]

    \item \textbf{(Fiedler vector concentration)}: For $\alpha \geq \alpha_{\text{crit}}$, the Fiedler vector $\mathbf{f}_2$ satisfies
    \[
        \max_{u \in S_1}\,f_{2,u} \;\geq\; \frac{1}{\sqrt{2}} \quad \text{and} \quad \min_{v \notin S_1}\,f_{2,v} \;\leq\; -\frac{1}{\sqrt{2n}}
    \]
    up to sign, meaning the vector concentrates its positive mass on the largest minority component $S_1$ and separates it from the rest.

    \item \textbf{(Near-optimality of Fiedler-aligned cuts)}: Any balanced partition $P$ whose boundary agrees with the sign change of $\mathbf{f}_2$ satisfies
    \[
        \mathrm{Obj}(P) \;\leq\; \mathrm{Obj}^* \cdot \left(1 + \frac{C_1}{\alpha}\right)
    \]
    for a constant $C_1$ depending only on $|E_r|$ and $k$.
\end{enumerate}
\end{proposition}

\begin{proof}[Proof Sketch]
\textbf{Part 1.} By the variational characterization,
\[
    \lambda_2(L_w) = \min_{\mathbf{x} \perp \mathbf{1},\, \mathbf{x}\neq 0} \frac{\mathbf{x}^\top L_w\, \mathbf{x}}{\|\mathbf{x}\|^2}
    = \min_{\mathbf{x} \perp \mathbf{1}} \frac{\displaystyle\sum_{(u,v)\in E} w(u,v)(x_u - x_v)^2}{\|\mathbf{x}\|^2}.
\]
Since $L_w = L_r + (\alpha-1)L_m$ and $L_m \succeq 0$,
\[
    \lambda_2(L_w) \;\geq\; \lambda_2(L_r) + (\alpha - 1)\,\lambda_2(L_m).
\]
If $G_m$ is connected ($c=1$), then $\lambda_2(L_m) > 0$ by the algebraic connectivity of $G_m$, giving $\lambda_2(L_w) \geq (\alpha-1)\lambda_2(L_m)$. In general, $\beta = \lambda_2(L_m) > 0$ whenever $G_m$ restricted to its largest component is connected, which holds whenever minority-dense tracts form a spatially contiguous region (as verified empirically for all five test states).

\textbf{Part 2.} For the upper bound on $\lambda_3$, note that when $G_m$ has $c \geq 2$ components, the constant vector supported on one component minus another component is orthogonal to $\mathbf{1}$ and to the Fiedler vector. Testing the Rayleigh quotient on the difference-of-indicators vector $\mathbf{z} = \mathbf{1}_{S_1}/\sqrt{|S_1|} - \mathbf{1}_{S_2}/\sqrt{|S_2|}$ shows $\mathbf{z}^\top L_m\, \mathbf{z} = 0$ (no $E_m$ edges cross components), so
\[
    \lambda_3(L_w) \;\leq\; \frac{\mathbf{z}^\top L_r\, \mathbf{z}}{\|\mathbf{z}\|^2} \;\leq\; \lambda_3(L_r),
\]
which is $O(1)$ in $\alpha$. The gap $\lambda_2/\lambda_3 = \Theta(\alpha)$ follows.

\textbf{Part 3.} Combine Davis--Kahan~\cite{davis1970rotation}: eigenvectors of $L_w$ change by $O(1/\lambda_2 - \lambda_3) = O(1/\alpha)$ under $O(1)$-magnitude perturbations. As $\alpha \to \infty$, $\mathbf{f}_2$ converges to the indicator vector of the largest minority component (normalized), giving the stated concentration bounds.

\textbf{Part 4.} Any partition respecting the sign pattern of $\mathbf{f}_2$ cuts only regular edges at the minority--non-minority boundary. The total cut cost is $\mathrm{Obj}^* + O(|E_r|/\alpha) = \mathrm{Obj}^*(1 + C_1/\alpha)$, since any deviation from the sign boundary to satisfy population balance cuts at most $O(|E_r|/k)$ additional regular edges at unit weight.
\end{proof}

\begin{remark}[Uniqueness of the Global Minimum]
\label{rem:uniqueness}
Proposition~\ref{prop:spectral} implies that for $\alpha \geq \alpha_{\text{crit}}$, the optimal balanced partition is essentially \emph{uniquely determined} by the minority subgraph structure of $G$: any balanced cut that avoids cutting minority--minority edges achieves near-optimal objective, and the set of such cuts is determined by the connected components of $G_m$. When those components are well-separated and there is a unique way to assign them to districts while satisfying the population balance constraint $\varepsilon$, the optimal partition is unique and every partitioning algorithm---regardless of tree structure---must find it. This is the spectral explanation for the empirical zero-variance finding.
\end{remark}

\subsection{Computational Implications}

Having characterized when method convergence occurs, we now analyze its computational implications.

\subsubsection{Time Complexity by Method}

Let $T_{\text{METIS}}(n, k)$ denote the runtime of a single METIS call partitioning $n$ vertices into $k$ parts. For well-separated clusters (as induced by $\alpha \geq \alpha_{\text{crit}}$), $T_{\text{METIS}}(n, k) = O(n \log k)$.

\begin{itemize}
    \item \textbf{Predetermined bisection}: One METIS call per level, $\log_2 k$ levels total.
    \[
    T_{\text{pred}} = O(\log k) \cdot T_{\text{METIS}}(n, 2) = O(n \log k)
    \]

    \item \textbf{Adaptive bisection}: At root level, evaluate $k$ candidate first splits. At subsequent levels, evaluate remaining splits.
    \[
    T_{\text{adapt}} = O(k) \cdot T_{\text{METIS}}(n, 2) + O(\log k) \cdot T_{\text{METIS}}(n, 2) = O(k \cdot n)
    \]

    \item \textbf{N-way partitioning}: Single call with $k$ parts.
    \[
    T_{\text{nway}} = T_{\text{METIS}}(n, k) = O(n k \log k)
    \]
\end{itemize}

\textbf{Observation}: For $\alpha \geq \alpha_{\text{crit}}$, adaptive incurs $O(k)$ overhead for zero benefit, since all candidate first splits produce identical downstream partitions.

\subsubsection{Early Termination for Adaptive Methods}

We propose an early termination criterion that eliminates adaptive overhead:

\begin{algorithm}[H]
\caption{Adaptive Bisection with Early Termination}
\begin{algorithmic}[1]
\STATE \textbf{Input:} Graph $G$, edge weights $w$, $k$ districts, tolerance $\epsilon$, threshold $\tau$
\STATE Evaluate first split candidate (e.g., $(1, k-1)$)
\STATE Check minority concentration: $c = \max_{i} M_i / N_i$
\IF{$c \geq \tau - \delta$ for some small $\delta$ (e.g., 0.05)}
    \STATE \textbf{return} this split (skip evaluating remaining $k-1$ candidates)
\ELSE
    \STATE Evaluate all $k$ candidate splits, select best
\ENDIF
\end{algorithmic}
\end{algorithm}

\textbf{Rationale}: For $\alpha \geq \alpha_{\text{crit}}$, the first split will achieve target concentration with high probability. Testing this criterion costs $O(1)$ but saves $O(k)$ METIS evaluations.

\subsubsection{Complexity-Theoretic Interpretation}

\begin{remark}[NP-Hardness Persists]
\label{rem:nphard}
A critical clarification is in order. \textbf{The balanced graph partitioning problem remains NP-hard even after introducing the edge weights $w$.} Specifically, for any fixed $\alpha \geq 1$, deciding whether a $k$-way balanced partition of a weighted graph has weighted cut value $\leq C$ is NP-complete by reduction from unweighted balanced bisection~\cite{garey1979}: replace each minority edge with an edge of weight $\alpha$ in the reduction instance.

What Theorem~\ref{thm:convergence} establishes is \emph{not} polynomial-time solvability, but rather a structural property of the \emph{solution landscape}: for $\alpha \geq \alpha_{\text{crit}}$, the weighted cut objective has a unique global minimum (under conditions C1--C3) and every balance-feasible local search algorithm---including METIS's multilevel Kernighan--Lin refinement---converges to a partition within $\delta(\alpha)$ of that minimum. The PFR pipeline with edge-weighting therefore finds a \emph{provably good local optimum} that happens to coincide with the global minimum, but it does so via a heuristic whose worst-case runtime is polynomial in $n$, not via an exact algorithm.

In formal terms: the edge-weighted partitioning problem on our VRA-redistricting instances is in the class of NP-hard problems whose \emph{smoothed complexity} is low~\cite{spielman2004smoothed}---random perturbations of any hard instance make it easy, and our geographic structure provides the structural analogue of such perturbations. This connects to the literature on ``easy instances of hard problems''~\cite{spielman2004}: worst-case hardness does not imply that typical instances arising from spatial demographic data are hard in practice.

Practitioners should therefore interpret the empirical equivalence result as follows: for the specific family of VRA-redistricting instances studied here, PFR with $\alpha = 5$ consistently finds the global optimum---but this is a consequence of the favorable structure of these instances, not a general guarantee extendable to arbitrary weighted graph partitioning inputs.
\end{remark}

\subsection{Summary of Theoretical Results}

Our theoretical framework establishes:

\begin{enumerate}
    \item \textbf{Method equivalence conditions} (Theorem~\ref{thm:convergence}): For $\alpha \geq \alpha_{\text{crit}}$ and under graph-structural conditions (C1)--(C3), all balance-feasible algorithms produce identical or near-identical partitions with $\mathrm{Obj}(P_{\mathcal{A}}) = \mathrm{Obj}^* + O(1/\alpha)$.

    \item \textbf{Phase transition} (Theorem~\ref{thm:phase-transition}): Sharp transition from method-dependent ($\alpha < \alpha_{\text{crit}}$) to method-independent ($\alpha \geq \alpha_{\text{crit}}$) regimes, with $\mathrm{Var}_{\mathcal{A}}[\mathrm{Obj}(P_{\mathcal{A}})]$ dropping from $\Theta(1)$ to $O(1/\alpha^2)$.

    \item \textbf{Bounds on $\alpha_{\text{crit}}$}: Approximately $k / m_{\text{state}}$, tightened by spatial clustering coefficient. For our test states, $\alpha_{\text{crit}} \in [3, 5]$; empirical ablation confirms zero method variance at $\alpha = 5$.

    \item \textbf{Spectral characterization} (Proposition~\ref{prop:spectral}, Remark~\ref{rem:uniqueness}): The eigenvalue gap $\lambda_2(L_w)/\lambda_3(L_w) = \Theta(\alpha)$ grows without bound as $\alpha \to \infty$, forcing the Fiedler vector to concentrate on the minority subgraph and making the globally optimal balanced cut unique and determined by minority component structure.

    \item \textbf{NP-hardness caveat} (Remark~\ref{rem:nphard}): The problem remains NP-hard even with edge weights. PFR + edge-weighting finds a provably good local optimum that coincides with the global minimum on these spatially structured instances, but this is an instance-specific property, not a general polynomial-time guarantee.

    \item \textbf{Computational implications}: Adaptive methods incur $O(k)$ overhead for zero benefit when $\alpha \geq \alpha_{\text{crit}}$; simple predetermined bisection is sufficient.
\end{enumerate}

These results collectively explain the empirical findings in Section~\ref{sec:results}: at $\alpha = 5 \geq \alpha_{\text{crit}} \approx 3$, the eigenvalue gap is $\lambda_2/\lambda_3 \approx 22.5$ (measured on Georgia), the Fiedler vector is highly stable under perturbation, and all three method families converge to the unique global minimum. The theory also provides a practitioner guideline: use $\alpha \geq 2k / m_{\text{state}}$ to ensure method independence while remaining well inside the phase-transition threshold.
